\section{Introduction}
\label{sec:introduction}

Large language model (LLM) agents face a fundamental constraint: the context
window is finite, but experience is not. An agent that converses across
hundreds of sessions accumulates knowledge, preferences, and relationships
that cannot fit in a single prompt. Memory---the ability to store, retrieve,
and \emph{manage} this accumulation over time---is now a first-class
architectural concern.

A companion survey~\cite{hui2026taxonomy} documents over 90 memory systems
that have emerged since 2023 and organizes them along four axes: structure,
retrieval, lifecycle, and preference emergence. That analysis reveals a
pronounced gap: the field has invested heavily in memory \emph{formation} and
\emph{retrieval} while largely neglecting memory \emph{management}. Only
18.2\% of surveyed systems implement consolidation (episodic-to-semantic
promotion), 12.5\% support transformation (reorganization, deduplication),
and 38.6\% implement any forgetting mechanism. Preference emergence---the
organic crystallization of behavioral patterns from accumulated
interaction---is present in only 4.5\% of systems.

This paper presents \alaya{}, an embeddable memory engine that addresses
these gaps by treating memory as a \emph{dynamic process} rather than a
static store. The design is grounded in two traditions that have studied
memory far longer than computer science:

\begin{enumerate}[leftmargin=*, nosep]
    \item \textbf{Computational neuroscience:} Complementary learning
    systems (CLS) theory~\cite{mcclelland1995} for dual-store consolidation,
    Bjork dual-strength forgetting~\cite{bjork1992} for principled decay,
    Hebbian learning~\cite{hebb1949} for use-dependent graph adaptation, and
    Collins--Loftus spreading activation~\cite{collins1975} for associative
    retrieval.

    \item \textbf{Yogac\=ara Buddhist psychology:} The theory of storehouse
    consciousness (\emph{\=alaya-vij\~n\=ana}, \cjk{阿賴耶識}) provides a
    process-oriented framework in which memories are \emph{seeds}
    (\emph{b\=ija}) that strengthen through \emph{perfuming}
    (\emph{v\=asan\=a}), ripen heterogeneously (\emph{vip\=aka}), and
    undergo periodic structural transformation
    (\emph{\=a\'sraya-par\=av\d{r}tti}). As we show in
    \Cref{sec:discussion}, these concepts converge with neuroscience on
    the same structural principles.
\end{enumerate}

\noindent \alaya{} translates these principles into a concrete system with
five lifecycle operations beyond storage and retrieval: consolidation,
transformation, forgetting, preference emergence, and emergent ontology
(unsupervised category formation). The system is implemented as a
zero-dependency Rust library with SQLite, requiring no external databases or
network services.

Our contributions are:
\begin{enumerate}[leftmargin=*, nosep]
    \item A three-store architecture (episodic, semantic, implicit) with a
    Hebbian graph overlay that reshapes through use, implementing the
    CLS/Yogac\=ara dual-system model.
    \item Five lifecycle operations that manage memory health over time:
    consolidation, transformation, forgetting, preference emergence, and
    emergent ontology.
    \item An emergent category formation mechanism that discovers semantic
    clusters through dual-signal analysis (embedding similarity + graph
    topology) without supervision.
    \item A zero-dependency, embeddable implementation in Rust with SQLite
    WAL and FTS5, available as an open-source crate.
    \item Baseline evaluation on LoCoMo and LongMemEval demonstrating a
    retrieval crossover that motivates lifecycle-aware memory management.
\end{enumerate}
